% SOURCE: Chinese PDF viewer pp. 140--141 (before heading 2).
\section{Conserved Quantities}
Symmetry is one of the most important properties of a system.

\textbf{Definition.} If the expectation value of a mechanical quantity does not change with time,
\[\frac{\dd\langle F\rangle}{\dd t}=0,\]
then the mechanical quantity $F$ is called a \emph{conserved quantity}.

From
\[\langle F\rangle=\mel{\psi}{\op F}{\psi}\]
and the Schr\"odinger equation
\[\ii\hbar\frac{\partial}{\partial t}\ket\psi=\op H\ket\psi,\]
we have
\begin{align*}
\frac{\dd\langle F\rangle}{\dd t}
&=\frac{\partial\bra\psi}{\partial t}\op F\ket\psi
 +\mel\psi{\frac{\partial\op F}{\partial t}}\psi
 +\bra\psi\op F\frac{\partial\ket\psi}{\partial t}\\
&=\left\langle\frac{\partial\op F}{\partial t}\right\rangle
 +\frac1{\ii\hbar}\langle[\op F,\op H]\rangle.
\end{align*}
If $\op F$ has no explicit dependence on time $t$,
\[\frac{\dd\langle F\rangle}{\dd t}=
\frac1{\ii\hbar}\langle[\op F,\op H]\rangle.\]
According to the definition, if $\op F$ and $\op H$ commute, then $F$ is a conserved quantity.

\textbf{Examples.}
\begin{enumerate}[label=\alph*)]
\item For a free-particle system, $\op H=\op{\vect p}^{\,2}/(2m)$; the quantity $\op{\vect p}$ does not depend explicitly on $t$ and $[\op{\vect p},\op H]=0$, so momentum is conserved.
\item For a general system, $\op H=\op{\vect p}^{\,2}/(2m)+V(\vect x)$, one generally has $[\op{\vect p},\op H]\ne0$; momentum is not conserved.
\item For a central-field system,
\[\op H=\frac{\op{\vect p}^{\,2}}{2m}+V(r)
=-\frac{\hbar^2}{2mr^2}\frac{\partial}{\partial r}
\left(r^2\frac{\partial}{\partial r}\right)
+\frac{\op{\vect L}^{\,2}}{2mr^2}+V(r),\]
orbital angular momentum $\op{\vect L}$ and its components do not depend explicitly on time, and $[\op{\vect L}^{\,2},\op H]=[\op{\vect L},\op H]=0$; orbital angular momentum and each of its components are conserved.
\item If $\op H$ has no explicit dependence on $t$, $[\op H,\op H]=0$, so energy is conserved.
\end{enumerate}
Whether a mechanical quantity is conserved is determined by the system's $\op H$.

\textbf{Properties of a conserved quantity.}
\begin{enumerate}[label=\alph*)]
\item Its expectation value in an arbitrary state is independent of time (by definition).
\item The probabilities of its possible measured values in an arbitrary state are independent of time.
\end{enumerate}
\textbf{Proof.} If $[\op F,\op H]=0$, the operators $\op F$ and $\op H$ have a common complete set of eigenstates $\ket n$,
\[\op F\ket n=F_n\ket n,\qquad \op H\ket n=E_n\ket n.\]
For an arbitrary state,
\[\ket{\psi(t)}=\sum_n\ket n\braket{n}{\psi(t)},\qquad
C_n(t)=\braket{n}{\psi(t)}.\]
The probability for a measurement of $\op F$ to give the value $F_n$ is $|C_n(t)|^2$. Since
\begin{align*}
\frac{\dd C_n(t)}{\dd t}
&=\bra n\frac{\partial}{\partial t}\ket{\psi(t)}
=\frac1{\ii\hbar}\mel n{\op H}{\psi(t)}\\
&=\frac{E_n}{\ii\hbar}\braket n{\psi(t)}
=\frac{E_n}{\ii\hbar}C_n(t),
\end{align*}
it follows that
\[C_n(t)=C_n(0)e^{-\ii E_nt/\hbar},\qquad
|C_n(t)|^2=|C_n(0)|^2,\]
independent of time.

\textbf{Inferences.}
\begin{enumerate}[label=\alph*)]
\item If the system initially occupies an eigenstate of a conserved quantity, it remains in that eigenstate.
\item If the system initially does not occupy an eigenstate of a conserved quantity, it never occupies an eigenstate of that conserved quantity.
\item In quantum mechanics, the quantum numbers of the eigenvalues of mechanical quantities are customarily used to label states; for example, the state $\ket{nlm}$ in a central field is described using the quantum numbers of energy, angular momentum and its component. However, the quantum numbers of a nonconserved quantity are unsuitable for describing a state: even if the initial state is an eigenstate of these quantities and can be described by their quantum numbers, the evolved state is no longer an eigenstate of those quantities, and those quantum numbers cannot be used to describe it. Only the quantum numbers of conserved quantities are good quantum numbers for describing states: if the system initially occupies an eigenstate of a conserved quantity, it remains in that eigenstate.
\end{enumerate}

\textbf{Distinguish two concepts.}

\emph{Conserved quantity:} Both the expectation value and the measurement probabilities of the conserved quantity in an arbitrary state do not vary with time.

\emph{Stationary state:} The expectation values and measurement probabilities of any mechanical quantity in a stationary state do not vary with time.